\chapter{Computable Numbers and Functions}

\section{Interval Arithmetic and Real Numbers \humanchecked}

\begin{definition}[Rational interval]\label{def:qinterval}
\lean{ComputableAnalysis.QInterval}
A \emph{rational interval} is a pair of rational numbers
\[
  [a,b]\in\Q\times\Q,\qquad a\le b.
\]
This is a formal interval object, not the set of rational points between
\(a\) and \(b\).  The \emph{width} of the rational interval is defined as \(b-a\).

Two rational intervals \([a, b]\) and \([c, d]\) \emph{overlap} when
\[
  a\le d\ \text{and}\ c\le b.
\]
\end{definition}

\begin{definition}[Interval sequence]\label{def:interval-sequence}
\uses{def:qinterval}
An \emph{interval sequence} is a sequence of rational intervals
\[
  [a_n,b_n]_{n\in\N},
  \qquad a_n \leq b_n\in\Q.
\]
\end{definition}

\begin{definition}[Real number]\label{def:raw-real}
\lean{ComputableAnalysis.RealRaw}
\lean{ComputableAnalysis.RealRaw.PrecisionValid}
\lean{ComputableAnalysis.RealRaw.PrecisionValidCompute}
\uses{def:interval-sequence}
A \emph{real number} is an interval sequence \(x=[a_n,b_n]_{n\in\N}\) satisfying the
following validity conditions:
\[
  n\le m\Rightarrow
  a_n\le a_m\le b_m\le b_n,
\]
and for every positive stage \(n\),
\[
  b_n-a_n < \frac{1}{n}.
\]
Stage \(0\) is allowed to be a convenient initial interval; the quantitative
precision contract starts only at \(n>0\).
\end{definition}

\begin{remark}[Computable number]\label{rem:computable-presentations}
\uses{def:interval-sequence,def:raw-real}
When the map 
\[
  n\mapsto [a_n,b_n]
\]
is given by a finite, terminating algorithm, as we always do in this project,
the real number is said to be \emph{computable}. Storing all the rational numbers
in an infinite table is not permitted. Evaluating at \(n\) is referred to as
\emph{computing} the real number \emph{at stage} \(n\).
\end{remark}

\begin{definition}[Positive and negative real numbers]\label{def:raw-real-sign}
\lean{ComputableAnalysis.RealRaw.Pos}
\lean{ComputableAnalysis.RealRaw.Neg}
\uses{def:raw-real}
Let \(x=[a_n,b_n]_{n\in\N}\) be a real number.  We say that \(x\) is
\emph{positive} (resp. \emph{negative}) if there is some \(N\in\N\) such that
\( a_N > 0\) (resp. \( b_N<0 \)).
The number \(x\) is \emph{nonzero} if it is either positive or negative.
\end{definition}

\begin{definition}[Interval Arithmetic]\label{def:raw-real-arithmetic}
\lean{ComputableAnalysis.RealRaw.add}
\lean{ComputableAnalysis.RealRaw.sub}
\lean{ComputableAnalysis.RealRaw.scaleRat}
\lean{ComputableAnalysis.RealRaw.scaleRatStage}
\lean{ComputableAnalysis.RealRaw.mul}
\lean{ComputableAnalysis.RealRaw.mulStage}
\lean{ComputableAnalysis.RealRaw.mulPrecisionFactor}
\uses{def:raw-real,def:raw-real-sign}
Let \(x\) and \(y\) be real numbers:
\[
  x=[a_n,b_n]_{n\in\N},\qquad y=[c_n,d_n]_{n\in\N}.
\]
The arithmetic operations are algorithms with their own internal precision
schedules.  For addition, a request for stage \(n\) asks each input for stage
\(2n\):
\[
  (x+y)_n := [a_{2n}+c_{2n},b_{2n}+d_{2n}],
\]
Negation of \([c_n,d_n]\) is given by \([-d_n,-c_n]\).  Hence subtraction uses
the same doubled request:
\[
  (x-y)_n := [a_{2n}-d_{2n},b_{2n}-c_{2n}].
\]
For a rational scalar \(r\), the formal algorithm first chooses a linear
internal stage depending on \(|r|\), then scales the resulting interval.  The
point of these stage choices is that a requested output precision can be
passed directly to the operation; the operation itself manages the precision
requested from its subroutines.

For products, first take rational bounds from the initial intervals.  For
example set
\[
  A_x=1+\max(|a_0|,|b_0|),\qquad
  A_y=1+\max(|c_0|,|d_0|).
\]
Choose a natural number \(M(x,y)\) large enough that
\[
  M(x,y) > A_x+A_y.
\]
The Lean definition implements this as a named stage factor computed from the
stage-zero rational endpoints; the proof only uses that it is a sufficiently
large positive integer.
At requested stage \(n\), compute both inputs at the internal stage
\[
  k=M(x,y)n
\]
and take the interval hull of the four endpoint products:
\[
  (xy)_n :=
  \left[
    \min S_k,\max S_k
  \right],
  \qquad
  S_k=\{a_kc_k,a_kd_k,b_kc_k,b_kd_k\}.
\]
The proof uses the product rule estimate
\[
  \operatorname{width}((xy)_n)
  \le A_x(d_k-c_k)+A_y(b_k-a_k)
\]
for \(k=M(x,y)n\).  Increasing \(M(x,y)\) by a fixed safety factor gives the strict bound
\(\operatorname{width}((xy)_n)<1/n\) for every \(n>0\).

If \(x\) is nonzero, let \(N\) be the smallest natural number such that
\([a_N,b_N]\) avoids zero.  Then
we define
\[
  x^{-1}:= [ b_{N+n}^{-1}, a_{N+n}^{-1} ]_{n\in\N}.
\]
Then \(y/x\) is defined as \(yx^{-1}\), when \(x\) is nonzero.
\end{definition}

\begin{definition}[Square-root subroutine]\label{def:sqrt-subroutine}
\lean{ComputableAnalysis.sqrtPartialRaw}
\lean{ComputableAnalysis.sqrtApproxOnDomain}
\lean{ComputableAnalysis.sqrtInitialWidthBound}
\lean{ComputableAnalysis.dyadicFuel}
\lean{ComputableAnalysis.dyadicFuel_bound}
\lean{ComputableAnalysis.sqrtStageFuel}
\uses{def:qinterval,def:raw-real-arithmetic}
For \(q\in\Q_{\ge 0}\), the square-root algorithm starts from the rational
interval
\[
  [0,\max(1,q)]
\]
and repeatedly bisects it.  At each step, if the midpoint \(m\) satisfies
\(m^2\le q\), replace the lower endpoint by \(m\); otherwise replace the upper
endpoint by \(m\).  Put \(U=\max(1,q)\), choose the integer
\[
B(q)=|\operatorname{num}(U)|+1,
\]
so that \(U<B(q)\).  Define the dyadic fuel \(D(T)\) recursively by halving
the target \(T\) until it is at most \(1\), with the proved bound
\[
  T<2^{D(T)}.
\]
At stage \(n\), run
\[
  D((n+1)B(q))
\]
bisection steps.  After \(t\) bisections the width is exactly \(U/2^t\).  Since
\((n+1)B(q)>nU\) and \((n+1)B(q)<2^{D((n+1)B(q))}\), this returns an
interval \([\ell_n,u_n]\) such that, for every \(n>0\),
\[
  0\le \ell_n\le u_n,\qquad
  \ell_n^2\le q\le u_n^2,\qquad
  u_n-\ell_n<\frac1n.
\]
This is the primitive used for rational segment lengths in the circumference
algorithm.  When a larger construction sums many square-root subroutines, it
must request a sufficiently large internal precision from each square-root
call, just as addition, scaling, and multiplication do for their inputs.
\end{definition}


\begin{proposition}[Arithmetic preserves real numbers]\label{prop:raw-arithmetic-valid}
\lean{ComputableAnalysis.RealRaw.add_precisionValid}
\lean{ComputableAnalysis.RealRaw.sub_precisionValid}
\lean{ComputableAnalysis.sqrtRaw_valid}
\lean{ComputableAnalysis.sqrtRaw_spec}
\uses{def:raw-real,def:raw-real-arithmetic}
The operations \(x+y\), \(x-y\), \(-x\), and \(xy\) are real numbers
whenever \(x\) and \(y\) are.  If \(x\) is nonzero, then \(x^{-1}\)
is also a real number.  For every rational \(q\ge 0\), the square-root
subroutine defines a valid raw real satisfying the square-root enclosure
invariant at every stage.
\end{proposition}

\begin{definition}[Equivalence of real numbers]\label{def:raw-equiv}
\lean{ComputableAnalysis.RealRaw.Equiv}
\uses{def:raw-real}
Two real numbers
\[
  x=[a_n,b_n]_{n\in\N},\qquad y=[c_n,d_n]_{n\in\N}
\]
are \emph{equivalent} if they overlap at every stage, i.e.,
\[
  [a_n,b_n] \mbox{ overlaps with } [c_n,d_n]
  \qquad\text{for all }n.
\]
\end{definition}

\begin{proposition}[Equivalence relation]\label{prop:raw-equiv-relation}
\lean{ComputableAnalysis.RealRaw.equiv_refl}
\lean{ComputableAnalysis.RealRaw.equiv_symm}
\lean{ComputableAnalysis.RealRaw.equiv_trans}
\uses{def:raw-real,def:raw-equiv}
Equivalence of real numbers is an equivalence relation.
\end{proposition}

\begin{definition}[Stage schedules]\label{def:stage-schedule}
\lean{ComputableAnalysis.RealRaw.StageSchedule}
\lean{ComputableAnalysis.RealRaw.schedule}
\uses{def:raw-real,def:raw-equiv}
A \emph{stage schedule} is an increasing function
\[
  \sigma:\N\to\N.
\]
If \(x=[a_n,b_n]_{n\in\N}\), then the scheduled real number is
\[
  \sigma\ldot x:=[a_{\sigma(n)},b_{\sigma(n)}]_{n\in\N}.
\]
\end{definition}

\begin{proposition}[Using stage schedules]\label{prop:stage-schedule-equiv}
\lean{ComputableAnalysis.RealRaw.schedule_equiv}
\lean{ComputableAnalysis.RealRaw.equiv_of_schedule_equiv}
\uses{def:stage-schedule,def:raw-equiv}
For every stage schedule \(\sigma\), one has
\[
  x\sim \sigma\ldot x.
\]
Thus, to prove \(x\sim y\), it is enough to choose schedules
\(\sigma,\tau\) that make
\[
  \sigma\ldot x\sim \tau\ldot y
\]
easier to show.
\end{proposition}

\begin{definition}[Abstract real numbers]\label{def:reals-quotient}
\lean{ComputableAnalysis.Real}
\uses{def:raw-real,def:raw-equiv}
The \emph{abstract real numbers} are the equivalence classes of real numbers
under the equivalence relation above.  We denote this quotient by \(\R\):
\[
  \R:=\{\text{real numbers}\}/\sim.
\]
After passing to the quotient, a real number in the sense above is a
\emph{representative} of the abstract real number it determines.
\end{definition}

\begin{theorem}[\(\R\) extends \(\Q\) as an ordered field]\label{thm:raw-reals-field}
\lean{ComputableAnalysis.Real.ofRat}
\lean{ComputableAnalysis.RealRaw.add_equiv}
\lean{ComputableAnalysis.RealRaw.sub_equiv}
\uses{def:reals-quotient,def:raw-real-sign,def:raw-real-arithmetic,prop:raw-arithmetic-valid}
The arithmetic operations and order on real numbers are compatible with
equivalence, and therefore descend to equivalence classes.  With these
descended operations and order, \(\R\) is an ordered field.
The map \(\Q\hookrightarrow\R\) that sends \(q\) to the class of
\( [q,q]_{n\in\N} \)
is a (homomorphic) embedding of ordered fields.
\end{theorem}

\begin{remark}[Completeness]\label{rem:real-completeness}
\uses{thm:raw-reals-field}
The definition above is equivalent to the standard definition by equivalence
classes of Cauchy sequences of rational numbers, hence \(\R\) is complete.
In this project, however, we avoid using completeness as a black box by
working with computable numbers as much as possible.  When a proof needs
existence or convergence, we build the relevant interval sequence and prove
its validity directly.
\end{remark}

\section{Examples of real numbers \humanchecked}

\begin{remark}[Types of real numbers]\label{rem:sources-of-raw-reals}
\uses{def:interval-sequence,def:raw-real}
In practice, there are various ways of defining real numbers, most of which
fall into one of the following types:
\begin{itemize}
  \item Algebraic numbers, i.e., solutions to polynomial equations such as
  \(x^2-2=0\).  Newton's method, or tangent-secant method, works well here, starting with a rational
  interval such as \([1,2]\).

  \item Infinite series, for example,
  \[
    \frac{\pi}{4}
    =
    1-\frac{1}{3}+\frac{1}{5}-\frac{1}{7}+\cdots
  \]
  and
  \[
    \frac{\pi^2}{6}
    =
    1+\frac{1}{4}+\frac{1}{9}+\frac{1}{16}+\cdots .
  \]
  of Leibniz and Euler. These are easy to implement, though the latter requires a little
  care with the upper bound (hint: telescoping sum).

  \item Continued fractions, such as
  \[
    \frac{4}{\pi}
    =
    1+\cfrac{1^2}{2+\cfrac{3^2}{2+\cfrac{5^2}{2+\cfrac{7^2}{2+\ddots}}}}.
  \]
  and finite truncations work as rational bounds.

  \item Periods, i.e., integrals of rational or algebraic functions:
  \[
    \frac{\pi}{4}
    =
    \int_0^1\frac{\dd x}{1+x^2},
  \]
  \[
    \frac{\pi}{6}
    =
    \int_0^{1/2}\frac{\dd x}{\sqrt{1-x^2}},
  \]
  \[
    \frac{\pi}{12}+\frac{\sqrt3}{8}
    =
    \int_0^{1/2}\sqrt{1-x^2}\,\dd x.
  \]
  and upper and lower Riemann sums work. (The last one is what Newton used to compute \(\pi\),
  with binomial expansion.)

  \item Integrals (proper or not) of transcendental functions are also
  likely in play, for instance
  \[
    \sqrt{2\pi}
    =
    \int_{-\infty}^{\infty} e^{-x^2/2}\,\dd x.
  \]
\end{itemize}

One may define \(\pi\) by one of these formulas, and the others would be
theorems (of equivalence of two real numbers). Or perhaps one should
define \(\pi\) by Archimedes' method of exhaustion.
\end{remark}

\begin{remark}[Calculus as equivalence]\label{rem:calculus-as-equivalence}
\uses{def:raw-equiv,rem:sources-of-raw-reals}
A great deal of calculus is the art of proving that independently constructed
real numbers are equivalent. The great advance by Newton and Leibniz is
to turn that art mechanical, by leveraging the fundamental formula
\[
  \int_a^b f'(x)\,\dd x = f(b)-f(a)
\]
which we shall establish, \emph{without} the completeness of real numbers.
\end{remark}

\begin{remark}[Algorithm with a variable]\label{rem:parameterized-constructions}
\uses{rem:sources-of-raw-reals}
In many of these algorithms, there is a natural variable that could assume
different values (\(\sqrt{\phantom{x}}\), \(\arctan\), \(\arcsin\) in the
examples above), so it makes sense to define the algorithms
\[
  n,t\longmapsto [a_n(t),b_n(t)]
\]
to be (computable) \emph{functions}. However, not all rational values of \(t\)
need to result in a real number.

This turns out to be crucial for establishing many of the equivalences of real
numbers, and is how calculus appears to be the study of functions.
\end{remark}

\begin{remark}[Complex variable]\label{rem:complex-parameters}
\uses{rem:parameterized-constructions,def:qcomplex,def:qbox}
It turns out that often it makes sense for the variable to assume complex values
(at least for \(t+si\), with \(t,s\in\Q\)).  For example, power series may be evaluated
inside a disk in the complex plane, and integrals can be taken over paths in the complex plane.
\end{remark}

\section{Complex numbers \humanchecked}

One could define complex numbers as pairs of real numbers.  Instead, we follow
the interval arithmetic more closely, using rectangular boxes with rational corners
from the start.

\begin{definition}[Field \(\Q(i)\)]\label{def:qcomplex}
\lean{ComputableAnalysis.QComplex}
Let
\[
  \Q(i):=\{r+is:r,s\in\Q\},
\]
with the standard field operations:
\[
  (r+is)(u+iv)=(ru-sv)+i(rv+su),
\]
etc. We also use the coordinatewise (partial) order:
\[
  r+is \le u+iv
  \quad\Longleftrightarrow\quad
  r\le u\ \text{and}\ s\le v.
\]
Note that this is not a total/linear order on \(\Q(i)\).
\end{definition}

\begin{definition}[Rational box]\label{def:qbox}
\lean{ComputableAnalysis.QBox}
\uses{def:qcomplex}
A \emph{rational box} is a pair
\[
  \cbox{a}{b}\in \Q(i)\times\Q(i),
  \qquad a\le b.
\]
Two boxes \(\cbox{a}{b}\) and
\(\cbox{c}{d}\) \emph{overlap} if
\[
  a\le d\ \text{and}\ c\le b.
\]
\end{definition}

\begin{definition}[Box sequence]\label{def:box-sequence}
\uses{def:qbox}
A \emph{box sequence} is a sequence of rational boxes
\[
  \cbox{a_n}{b_n}_{n\in\N}.
\]
\end{definition}

\begin{definition}[complex number]\label{def:complex-raw}
\lean{ComputableAnalysis.ComplexRaw}
\uses{def:box-sequence,def:raw-equiv}
A \emph{complex number} is a nested box sequence
\[
  z=\cbox{a_n}{b_n}_{n\in\N}
\]
satisfying
\[
  n\le m\Rightarrow a_n\le a_m\le b_m\le b_n,
\]
and whose real and imaginary widths both shrink to \(0\).  Equivalently, its
coordinate interval sequences are real numbers.  Two complex numbers are
equivalent when their boxes overlap at every stage, i.e. coordinatewise
equivalence of the real and imaginary parts.
\end{definition}

\begin{definition}[Complex arithmetic]\label{def:complex-arithmetic}
\lean{ComputableAnalysis.ComplexRaw.add}
\lean{ComputableAnalysis.ComplexRaw.mul}
\uses{def:complex-raw,def:raw-real-arithmetic}
For complex numbers as box sequences
\[
  z=\cbox{a_n}{b_n}_{n\in\N},\qquad
  w=\cbox{c_n}{d_n}_{n\in\N},
\]
addition is stagewise:
\[
  z+w:=\cbox{a_n+c_n}{b_n+d_n}_{n\in\N}.
\]
Negation reverses the box:
\[
  -z:=\cbox{-b_n}{-a_n}_{n\in\N}.
\]
For multiplication, write the same box sequences in coordinates as
\[
  z=x+iy,\qquad w=u+iv,
\]
where \(x,y,u,v\) are real numbers.  Then define
\[
  zw:=(xu-yv)+i(xv+yu).
\]
At the level of rational boxes, multiplication is implemented by applying
real interval arithmetic to these four real products and taking bounding
intervals.
\end{definition}

\begin{definition}[Complex inverse]\label{def:complex-inverse}
\lean{ComputableAnalysis.ComplexRaw.inv}
\uses{def:complex-arithmetic,def:raw-real-arithmetic,def:raw-real-sign}
A complex number \(z=x+iy\) is \emph{nonzero} if
\[
  x^2+y^2
\]
is nonzero as a real number.  For nonzero \(z\), define
\[
  z^{-1}
  :=
  \frac{x}{x^2+y^2}
  -
  i\,\frac{y}{x^2+y^2}.
\]
\end{definition}

\begin{definition}[Abstract complex numbers]\label{def:complex-quotient}
\lean{ComputableAnalysis.Complex}
\uses{def:complex-raw}
The \emph{abstract complex numbers} are the equivalence classes of complex
numbers under the equivalence relation above.  We denote this quotient by
\(\C\):
\[
  \C:=\{\text{complex numbers}\}/\sim.
\]
After passing to the quotient, a complex number in the sense above is a
\emph{representative} of the abstract complex number it determines.
\end{definition}

\begin{theorem}[\(\C\) is a field]\label{thm:complex-field}
\uses{def:complex-quotient,def:complex-arithmetic,def:complex-inverse}
The arithmetic operations above descend to abstract complex numbers, and
\(\C\) forms a field.
\end{theorem}

\begin{theorem}[\(\C\) extends \(\Q(i)\) and \(\R\)]\label{thm:rational-complex-field-square}
\lean{ComputableAnalysis.ComplexRaw.ofQComplex}
\lean{ComputableAnalysis.Complex.ofQComplex}
\lean{ComputableAnalysis.Complex.ofReal}
\uses{def:qcomplex,def:complex-quotient,thm:raw-reals-field,thm:complex-field}
The natural embeddings of fields form a commutative square
\[
\begin{array}{ccc}
\Q & \longrightarrow & \Q(i)\\
\downarrow & & \downarrow\\
\R & \longrightarrow & \C
\end{array}
\]
of field homomorphisms.
\end{theorem}

\section{Functions and domains}

\begin{definition}[Complex-valued function representation]\label{def:function-raw}
\lean{ComputableAnalysis.FunctionRaw}
\uses{def:qcomplex,def:complex-raw}
A \emph{complex-valued function representation} consists of a rational domain
predicate
\[
  D_f\subseteq \Q(i)
\]
and, for every \(q\in D_f\), a complex number
\[
  f(q)=\bigl(B_n(q)\bigr)_{n\in\N}.
\]
Validity says that each pointwise box sequence \(f(q)\) is a complex number.
This is the underlying notion even when a concrete function is intended only
for real inputs or real outputs.
\end{definition}

\begin{definition}[Real-input and real-valued functions]\label{def:real-fun-raw}
\lean{ComputableAnalysis.RealFunRaw}
\uses{def:function-raw,thm:rational-complex-field-square}
A real-input function is a complex-valued function representation whose domain
is restricted to the rational real axis \(q=x+i0\).  It is \emph{real-valued}
when each output box lies on the real axis, equivalently when its imaginary
coordinate is the zero real number.

The current real-function layer records this common special case directly as
a rational domain predicate \(D_f\subseteq\Q\) and rational interval outputs.
It should be viewed as a convenient real-axis presentation of the complex
function layer, not as a separate foundation.
\end{definition}

\begin{definition}[Rational segment]\label{def:rational-segment}
\uses{def:qcomplex}
A \emph{rational segment} in the complex plane is determined by two distinct
rational complex points \(p,q\in\Q(i)\).  Its rational points are
\[
  (1-t)p+tq,\qquad t\in\Q,\quad 0\le t\le 1.
\]
The segment is real when \(p\) and \(q\) lie on the rational real axis.
\end{definition}

\begin{definition}[Agreement on an overlap]\label{def:function-agrees-on-overlap}
\uses{def:function-raw,def:rational-segment,def:complex-raw}
Let \(f\) and \(g\) be complex-valued function representations.  An
\emph{agreement-on-overlap certificate} consists of a nontrivial rational
segment \(S\) such that every rational point of \(S\) lies in both domains, and
for every rational point \(q\in S\), the complex numbers \(f(q)\) and \(g(q)\)
are equivalent.

This is a local agreement statement with an explicit witness.  For arbitrary
functions it is not promoted to a global equivalence relation; transitivity
requires additional analytic structure.
\end{definition}

\begin{definition}[Function on an interval]\label{def:function-on-interval}
\lean{ComputableAnalysis.FunctionOnInterval}
\uses{def:qinterval,def:real-fun-raw,def:rational-segment}
A function on a rational interval \([a,b]\) is the real-axis special case of
the preceding setup: the interval is the rational segment from \(a+i0\) to
\(b+i0\), the function is defined at every rational point of that segment, and
the pointwise outputs are real numbers.
\end{definition}

\begin{definition}[Interval regularity]\label{def:interval-regular}
\lean{ComputableAnalysis.IntervalRegularOn}
\uses{def:function-on-interval}
Let \(F\) be a real-valued function on \([a,b]\).  An interval-regularity
certificate assigns to each rational subinterval
\[
  K=[r,s]\subseteq [a,b]
\]
and each stage \(n\) a rational output interval \(Y(K,n)\).  It supplies a
modulus \(M:\N\to\N\) such that whenever
\[
  s-r\le \frac{1}{M(n)},
\]
and \(Y(K,n)=[u,v]\), we have
\[
  0\le v-u\le \frac{1}{n},
\]
and for every rational point \(x\in K\), the point-value interval
\(F(x,n)\) overlaps \(Y(K,n)\).
\end{definition}

\begin{remark}[Holomorphic branch equivalence]\label{rem:holomorphic-branch-equivalence}
\uses{def:function-agrees-on-overlap}
For later complex analysis, holomorphic function representations will carry
stronger domain data.  On that layer, overlap should be witnessed by a
nonempty open region, such as a rational disk contained in the intersection of
the domains.  The identity principle can then turn local agreement on such an
overlap into a transitive branch-equivalence relation.
\end{remark}

\begin{remark}[Completeness is replaced by moduli]\label{rem:completeness-replaced-by-moduli}
\uses{def:interval-regular}
Classical compactness, uniform continuity, and existence of suprema are not
used as black boxes.  The project records the finite data that those theorems
usually hide: interval images, rates, nesting, and overlap proofs.
\end{remark}
